% flour-types-ar.tex — Arabic edition of ch. 7, flour types
%
% Original: The Sourdough Framework by Hendrik Kleinwächter,
% the-sourdough-framework.com, CC BY-SA 4.0. The Arabic edition inherits
% CC BY-SA 4.0.
%
% Phase 28R R2 unit 9 reconstruction. Body text recovered verbatim from
% the R0 recovery reference (recovery-r0/text/07-flour-types.txt,
% mined from the Phase 23 website build of the accepted Arabic sources)
% with the accepted edition book-ar-fixed.pdf (b78c2a28b574288801d8000a7249af3a)
% as the fidelity authority; no text retranslated from the English.
% Figure/table labels recovered from the accepted edition's renders
% (pp. 56–59).
%
% Direction treatments probe-verified 2026-11-15 (probe9 series):
% body-text T-numbers ride in per-token \LR islands (a bare T405
% renders "405T" in an RTL line); tables 7.1/7.2 are mirrored and
% table 7.3 is NOT mirrored, per the accepted edition's renders; mixed
% "or" cells are written in reversed literal order (probe9c); the
% الترطيب (%) header of table 7.3 carries unwrapped parentheses per
% the accepted render )%(. The (> 70%) cell of table 7.2 carries a
% documented paren-font residual (Pagella vs the accepted Amiri).

\chapter{أنواع الدقيق}\label{ch:flour-types}

\begin{quoting}
في هذا الفصل سنلقي نظرة فاحصة على أنواع الدقيق المختلفة وطريقة تصنيف كلٍّ منها. سنستعرض أيضًا الطرق الشائعة للتمييز بين أنواع الدقيق المتشابهة، لتشتري الدقيق الذي تحتاجه بثقة أكبر.
\end{quoting}

أبسط أنواع الدقيق هو دقيق كامل الحبة \textenglish{(whole grain flour)}، وفي هذه الحالة تكون الحبة بأكملها قد طُحنت إلى قطع أصغر. وأحيانًا، بحسب ما تريد أن تخبزه، قد لا تكون النكهة القوية للنخالة مرغوبة؛ وعندها يمكنك استخدام أنواع دقيق أكثر بياضًا. فالمطاحن تزيل، بمساعدة المناخل \textenglish{(sieves)}، الأجزاء الكبيرة من غلاف الحبّة \textenglish{(hull)}. وتحمل الحبة أصلًا جنينًا \textenglish{(germ)} جاهزًا من النبات ينتظر أن يُستثار. وأبيض دقيق يمكنك الحصول عليه هو في معظمه جزء النشا \textenglish{(starch)} من الحبة وحده. وحسب الطبقات الباقية في الدقيق تُستخدم أسماء مختلفة لوصف نوعه.

\begin{table}[!htb]
\centering
% table-flour-types-ar.tex — Arabic edition of upstream
% tables/table-flour-types.tex (ch. 7, جدول 7.1)
%
% Original: The Sourdough Framework by Hendrik Kleinwächter,
% the-sourdough-framework.com, CC BY-SA 4.0. The Arabic edition inherits
% CC BY-SA 4.0.
%
% RTL mirror per the settled table mechanics: the column order is
% reversed in the declared LTR tabular (upstream {@{}llrrr@{}} becomes
% {@{}rrrll@{}}) so إيطاليا renders rightmost and الولايات المتحدة
% leftmost, matching the accepted edition's render of p. 56 (table 7.1
% is mirrored; alignment follows the reversed columns: the three
% T-number columns right-aligned, the two text columns left-aligned).
%
% CELL DIRECTION (unit 9 finding, 2026-11-15): the accepted edition
% renders table cells as RTL paragraphs (first logical word rightmost:
% دقيق الكعك shows الكعك leftmost, دقيق rightmost). The sty's
% \insert@column patch appends \textdir TRT AFTER \@sharp — after the
% cell text — and is therefore a silent no-op; every multi-word Arabic
% cell in the delivered units 5–8 tables renders LTR word order
% against the accepted RTL (latent deviation, reported for the R2/R3
% rebuild). Until the sty ruling, unit 9 wraps every non-empty cell in
% {\textdir TRT ...} at the content layer (probe9g/probe9h verified
% the mechanism against the accepted orders, headers included).
%
% CONTENT NOTE — the accepted edition's Italy column diverges from
% upstream: rows 1--2 carry دقيق الكعك and متعدد الاستخدامات where
% upstream has the bare 00 and 0, and rows 4--5 are empty where
% upstream has 1 and 2. The accepted PDF (p. 56) and the R0 recovery
% reference agree line for line on this, so the accepted content is
% reproduced here (the accepted edition is the fidelity authority).
% This divergence also sets the Italy column's width (as wide as
% متعدد الاستخدامات), which the per-char battery confirmed against
% the accepted render.
%
% With TRT cells the mixed "or" cells carry the NATURAL logical order
% (T405 أو T550, 00 أو 0) and render exactly as the accepted edition
% (T550 أو T405 / 0 أو 00 in x-order) — the reversed-literal-order
% workaround of probe9c is superseded by the TRT cells. Bare digit
% runs ride on the lua's TRT digit handling (probe9g T5); bare
% T-numbers still need per-token \LR islands (the letter–digit swap,
% probe9b); Latin words ride in \LR islands per the unit 8 convention.

\begin{tabular}{@{}rrrll@{}}
\toprule
\thead{{\textdir TRT إيطاليا}} & \thead{{\textdir TRT فرنسا}} & \thead{{\textdir TRT ألمانيا}} & \thead{{\textdir TRT بريطانيا}} & \thead{{\textdir TRT الولايات المتحدة}} \\
\midrule
{\textdir TRT دقيق الكعك} & {\textdir TRT \LR{T45}} & {\textdir TRT \LR{T405}} & {\textdir TRT دقيق طري} & {\textdir TRT كعك} \\
{\textdir TRT متعدد الاستخدامات} & {\textdir TRT \LR{T55}} & {\textdir TRT \LR{T550}} & {\textdir TRT دقيق عادي} & {\textdir TRT متعدد الاستخدامات} \\
{\textdir TRT 00 أو 0} & {\textdir TRT \LR{T45} أو \LR{T55}} & {\textdir TRT \LR{T405} أو \LR{T550}} & {\textdir TRT دقيق الخبز} & {\textdir TRT دقيق الخبز} \\
 & {\textdir TRT \LR{T80}} & {\textdir TRT \LR{T812}} & & \\
 & {\textdir TRT \LR{T110}} & {\textdir TRT \LR{T1050}} & & \\
{\textdir TRT \LR{Integrale}} & {\textdir TRT \LR{T150}} & {\textdir TRT \LR{Vollkorn}} & {\textdir TRT كامل الحبة} & {\textdir TRT كامل الحبة} \\
\bottomrule
\end{tabular}

\caption[تسمية دقيق القمح]{مقارنة بين طريقة تسمية أنواع دقيق القمح المختلفة في دول مختلفة.}\label{tab:flour-types-comparison}
\end{table}

في ألمانيا يُستخدم محتوى الرماد \textenglish{(ash content)} لوصف أنواع الدقيق. يحرق المختبر \LR{\qty{100}{\gram}} من الدقيق في الفرن، ثم يُستخرج الرماد المتبقي ويُقاس. وحسب الكمية الناتجة يُصنَّف الدقيق. فإذا كان الدقيق من النوع \num{405}، فقد بقي بعد حرق الدقيق \LR{\qty{405}{\mg}} من الرماد. وكلما زادت أجزاء غلاف الحبّة في الدقيق بقيت معادن أكثر؛ لذلك كلما ارتفع الرقم اقترب الدقيق من الدقيق الكامل. وتختلف هذه الأرقام قليلًا بين كل نوع من الحبوب. وعلى العموم، كلما ارتفعت القيمة كان الطعم أقوى.

\begin{figure}[htb!]
\centering
\includegraphics[width=0.7\textwidth]{wheat-kernel-overview}
\caption[محتوى حبة القمح]{نظرة عامة على حبة القمح ومحتوياتها~\cite{wheat+kernel}.}\label{fig:wheat-kernel-overview}
\end{figure}

إذا قارنت بين أنواع الحبوب المختلفة ستجد حبوبًا عالية الغلوتين وأخرى منخفضة الغلوتين وأخرى خالية منه تمامًا. الغلوتين هو ما يمنح الخبز قوامه الإسفنجي. وبدون الغلوتين لن تمتلك المخبوزات الخصائص ذاتها. وإدارة الغلوتين تجعل عملية صناعة الخبز بأكملها أكثر تعقيدًا إذ تُضاف خطوات أكثر.

أما العجين الخالي من الغلوتين فلا يحتاج إلى عجن، لأن دور العجن هو تكوين روابط الغلوتين. وكلما زاد العجن قويت هذه الروابط. ومع الدقيق منخفض الغلوتين أو الخالي منه يكفي أن تخلط المكوّنات معًا، مع الحرص على تجنيسها \textenglish{(homogenize)} جميعها جيدًا.

أثناء التخمير تتفكك بروتينات الغلوتين إذ تستقلبها الكائنات الدقيقة. وحين يتحول قدر كبير من الغلوتين لن يحتفظ العجين بالبنية القمحية التي وصفناها سابقًا. وفي الدقيق الخالي من الغلوتين أو منخفضه يكون تركيزك الرئيس على إدارة الحموضة، فلا تريد أن يكون الخبز الناتج حامضًا أكثر من اللازم. وفي المقابل لن تقلقك مسألة تفكك الغلوتين، ما يزيل عبئًا كبيرًا من المعادلة.

\begin{table}[!htb]
\centering
% table-grains-bread-making-process-ar.tex — Arabic edition of upstream
% tables/table-grains-bread-making-process.tex (ch. 7, جدول 7.2)
%
% Original: The Sourdough Framework by Hendrik Kleinwächter,
% the-sourdough-framework.com, CC BY-SA 4.0. The Arabic edition inherits
% CC BY-SA 4.0.
%
% RTL mirror per the settled table mechanics: the column order is
% reversed in the declared LTR tabular (upstream {@{}lcccc@{}} becomes
% {@{}ccccc}) so نوع الحبّة renders rightmost and التشكيل leftmost,
% matching the accepted edition's render of p. 58 (table 7.2 is
% mirrored; all columns centered).
%
% SPEC NOTE (unit 9, 2026-11-15): the trailing @{} is dropped — the
% accepted edition keeps the leading @{} (التشكيل touches the left
% rule) but carries 6 pt of padding on the right edge (the widest
% الجاودار... cell ends 6 pt before the right rule), and the centered
% table then sits 3 pt left of the fully-@{}'d variant. Measured
% against accepted p. 58: rules 103.0--496.0, every cell width
% identical, all cells shifted exactly 3.0 pt left of the
% {@{}ccccc@{}} variant. probe9i S1--S4 ruled out the other spec
% variants (both @{}: 0 edge padding; no @{}: +12 pt; leading-only
% drop: padding on the wrong side).
%
% CELL DIRECTION: every non-empty cell wrapped in {\textdir TRT ...}
% at the content layer — the accepted edition renders cells as RTL
% paragraphs and the sty's \insert@column \textdir TRT patch is a
% silent no-op (placed after \@sharp); see table-flour-types-ar.tex
% for the full finding. probe9g/probe9h verified the mechanism.
%
% The (> 70%) cell rides in a single \LR island after the Arabic
% (island renders leftmost of the Arabic, matching the accepted cell
% of p. 58). Documented residual: the accepted edition renders the
% parentheses and > in Amiri with the digits and % in Pagella; every
% \LR variant available to the pipeline sets the whole island in
% Pagella (probe9e/probe9f, 2026-11-15, the \arabicfont wrap did not
% take), so the delivered cell carries Pagella parentheses — recorded
% as a paren-font residual, not a sequencing deviation.

\begin{tabular}{@{}ccccc}
\toprule
\thead{{\textdir TRT التشكيل}} & \thead{{\textdir TRT الطيّ والشدّ}} & \thead{{\textdir TRT العجن}} & \thead{{\textdir TRT التجنيس}} & \thead{{\textdir TRT نوع الحبّة}} \\
\midrule
{\textdir TRT نعم} & {\textdir TRT نعم} & {\textdir TRT نعم} & {\textdir TRT نعم} & {\textdir TRT السبلت والقمح \LR{(\textgreater{}~70\%)}} \\
{\textdir TRT لا} & {\textdir TRT لا} & {\textdir TRT لا} & {\textdir TRT نعم} & {\textdir TRT الجاودار والإممر والقمح الأحادي والأرز والذرة} \\
\bottomrule
\end{tabular}

\caption[أنواع الحبوب المختلفة]{نظرة عامة على أنواع الحبوب المختلفة والخطوات المتضمنة في عملية صناعة الخبز الخاصة بكلٍّ منها.}
\end{table}

لما كان للغلوتين دور خاص، فسيُخصَّص ما تبقى من هذا الفصل لنظرة فاحصة على أنواع الدقيق الغنية بالغلوتين وكيفية التمييز بينها. ومثل القمح، يحتوي السبلت على كميات كبيرة من الغلوتين، لذلك تنطبق عليه الخصائص ذاتها.

تستدعي وصفات عديدة دقيق الخبز القمحي \textenglish{(bread flour)}، غير أن مصطلح دقيق الخبز قد يشير إلى أنواع مختلفة من الدقيق. فقد يكون في ألمانيا من النوع \LR{T405} أو \LR{T550}، وكثيرًا ما يُساء تصنيف ذلك؛ فالمصطلحان الدقيق القوي \textenglish{(strong)} أو دقيق الخبز يشيران هنا إلى خواص الدقيق. فدقيق الخبز يُفهم على أنه دقيق يحتوي كمية أعلى من البروتين، ومن ثم من الغلوتين. وهذا الدقيق ممتاز حين تريد صنع الخبز الحامض لأن عجينتك تسمح بفترة تخمير أطول. وكما وصفنا سابقًا، يستهلك الكائنات الدقيقة الغلوتين. وكلما زاد الغلوتين احتفظت عجينتك بتماسكها مدة أطول. أما إذا أردت صنع كعكة فربما ترغب في دقيق أقل غلوتينًا، إذ قد لا تكون خواص تماسك الغلوتين مرغوبة لأن الكعكة الناتجة قد يغلب عليها قوام مطاطي لزج.

وخلاصة القول: ليس كل دقيق \LR{T405} أو \LR{T45} أو \LR{T00} مثل آخره. فحسب خواص النبات الذي جاء منه تختلف خواص الدقيق. ولهذا السبب أدخلت بعض الدول مثل ألمانيا سلالم إضافية لتقييم جودة القمح. فالفئة A تشير إلى قمح ذي جودة جيدة يمكن مزجه بجودات أدنى لتحسين الدقيق. والفئة B تشير إلى قمح متوسط يمكن استخدامه في صنع مخبوزات مختلفة. أما الفئة C فتخص قمحًا ضعيف الخصائص الخبازية، وقد يحدث ذلك مثلًا إذا كان القمح قد بدأ بالإنبات وبذلك فقد بعض خواصه الخبازية المرغوبة؛ وهذا النوع من القمح يُستخدم عادةً في علف الحيوانات أو ككتلة حيوية قابلة للتخمير في المولدات. والفئة E تشير إلى قمح النخبة \textenglish{(Elite)}، وهو أعلى جودة يمكن أن يبلغها القمح. وهذا النوع لا يُحصد إلا حين يكون القمح قد نما في ظل ظروف مثالية. ويمكنك مقارنة ذلك بمزرعة عنب تستخدم أفضل عنبها وحده لإنتاج نبيذ احتياطي. وللأسف لا تُطبع هذه الفئات عادةً على عبوة الدقيق الذي تشتريه. ويمكنك مراقبة قيمة البروتين بوصفها مؤشرًا محتملًا. غير أن المطاحن الكبيرة تمزج أنواع الدقيق معًا للحفاظ على الجودة على مدار السنين، ولا يُذكر الدقيق المزجي على العبوة أيضًا. وقد تستخرج المخابز الغلوتين من بعض أنواع الدقيق ثم تخلطه بهدف صنع دقيق خبز أفضل.

وفي إيطاليا طُوّر ما يُعرف بقيمة W \textenglish{(W-value)} لإظهار سلوك الدقيق على نحو أفضل. يُصنع عجين ثم يُقاس مدى مقاومة هذا العجين للعجن. فكلما زاد الغلوتين في الدقيق كان العجين أكثر مرونة، وكلما كان أكثر مرونة زادت مقاومتها للعجن. والدقيق ذو القيمة W الأعلى يحتوي نسبة غلوتين أعلى ويسمح بفترة تخمير أطول. لكن في المقابل يصعب على الكائنات الدقيقة نفخ العجين لأن كمية مادة البالون \textenglish{(balloon)} فيها أكبر. ولتصنع منتجًا مخمَّرًا ممتازًا من دقيق عالي القيمة W ستحتاج إلى فترة تخمير طويلة. والفترة الطويلة تعني أيضًا أن الكائنات الدقيقة ستُغني عجينتك بنكهة أكبر.

\begin{table}[!htb]
\centering
% table-overview-w-values-ar.tex — Arabic edition of upstream
% tables/table-overview-w-values.tex (ch. 7, جدول 7.3)
%
% Original: The Sourdough Framework by Hendrik Kleinwächter,
% the-sourdough-framework.com, CC BY-SA 4.0. The Arabic edition inherits
% CC BY-SA 4.0.
%
% NOT mirrored: the accepted edition renders this table in the upstream
% column order (p. 59: قيمة W leftmost, مدة التخمير rightmost), unlike
% tables 7.1 and 7.2 which are mirrored. The upstream spec
% {@{}rcll@{}} is kept as declared.
%
% CELL DIRECTION: every non-empty cell wrapped in {\textdir TRT ...}
% at the content layer — the accepted edition renders cells as RTL
% paragraphs and the sty's \insert@column \textdir TRT patch is a
% silent no-op (placed after \@sharp); see table-flour-types-ar.tex
% for the full finding. probe9g/probe9h verified the mechanism,
% headers included: the الترطيب (%) header renders )%(-mirrored
% parens leftmost of الترطيب exactly as the accepted p. 59 header,
% with the documented paren-font residual (accepted Amiri parens,
% pipeline Pagella).
%
% Ranges ride in \LR islands as LITERAL TEXT ranges —
% \LR{\textenglish{0\textendash 150}} etc. — per the unit 9 dash-font
% finding (2026-11-15): the accepted edition renders each range as a
% single Pagella token, and probe9l measured the literal text form at
% the accepted widths exactly (150–250 → 41.8 vs accepted 41.9;
% 0–150 / 50–60 / 60–70 / 70–90 → 29.9 vs 29.9) while \numrange with
% the \textenglish range-phrase came out 0.9–2.7 pt wider (siunitx
% digit spacing). The standing "--" → \textendash rule is preserved:
% the dash IS \textendash, wrapped in \textenglish so it takes Pagella
% rather than the Amiri a bare \textendash picks up inside a TRT cell
% (probe9j J1/J3). The قصيرة–متوسطة dash carries the same
% \textenglish{\textendash} wrapper (accepted p. 59, x 415.3–421.3,
% Pagella between the two Amiri words). \num{50} and \LR{350+} matched
% the accepted render as-is and stay.

\begin{tabular}{@{}rcll@{}}
\toprule
\thead{{\textdir TRT قيمة W}} & \thead{{\textdir TRT الترطيب (\%)}} & \thead{{\textdir TRT الاستخدامات}} & \thead{{\textdir TRT مدة التخمير}} \\
\midrule
{\textdir TRT \LR{\textenglish{0\textendash 150}}} & {\textdir TRT \LR{\num{50}}} & {\textdir TRT البسكويت} & {\textdir TRT قصيرة جدًا} \\
{\textdir TRT \LR{\textenglish{150\textendash 250}}} & {\textdir TRT \LR{\textenglish{50\textendash 60}}} & {\textdir TRT الكعك والخبز والبيتزا} & {\textdir TRT قصيرة\textenglish{\textendash}متوسطة} \\
{\textdir TRT \LR{\textenglish{250\textendash 350}}} & {\textdir TRT \LR{\textenglish{60\textendash 70}}} & {\textdir TRT الخبز والبيتزا} & {\textdir TRT طويلة} \\
{\textdir TRT \LR{350+}} & {\textdir TRT \LR{\textenglish{70\textendash 90}}} & {\textdir TRT الخبز والبيتزا} & {\textdir TRT طويلة جدًا} \\
\bottomrule
\end{tabular}

\caption[مدة التخمير مقابل قيمة W]{نظرة عامة على مستويات مختلفة من قيمة W وما يقابلها من ترطيب ومدد تخمير.}\label{tab:w-value}
\end{table}

وعلى العموم، حين تسعى لخبز رغيف قائم بذاته استهدف محتوى بروتين أعلى. فإذا كانت قيمة الغلوتين منخفضة نسبيًا سينهار خبزك أسرع. يبقى خبز الرغيف ممكنًا، لكن قد يكون أسهل أن تلجأ إلى تقنيات أخرى مثل قالب الخبز، أو أن تفكر في خبز المقلاة \textenglish{(skillet bread)} أو الخبز المُسطَّح \textenglish{(flatbread)}.

ومن الخصائص الإضافية النادرة الذكر لدقيق جيد درجة تضرر جزيئات النشا. وهذه مشكلة شائعة حين تحاول أن تطحن دقيق القمح بنفسك في المنزل؛ فالأرجح أن مطحنتك المنزلية غير قادرة على بلوغ النتائج نفسها التي يبلغها مطحن أكبر. وتضرر النشا أمر أساسي لتحسين خواص العجين. فستحصل على جيلاتنة \textenglish{(gelatinization)} وامتصاص ماء أفضل مع نشا متضرر كما ينبغي~\cite{starch+damage+flour}. فكلما زاد النشا المتضرر اتسعت المساحة السطحية، وهذا يحسّن طريقة تفاعل الماء مع الدقيق. كما يوفر سطحًا أكبر تستخدمه الكائنات الدقيقة لمهاجمة الجزيئات وبدء عملية التخمير.

ما زلت حتى الآن بلا طريقة جيدة لطحن دقيق القمح في المنزل. فحتى بعد أن حاولت طحن الدقيق \num{10} مرات مع فترات راحة قصيرة، لم أستطع بلوغ الخواص نفسها التي يبلغها الدقيق المطحون تجاريًا. كانت العجائن التي أعدتها تبدو جيدة، وربما بملمس خشن قليلًا. لكن أثناء الخَبز كانت تبدأ بفقد غازاتها سريعًا وتتحول إلى خبز مسطّح جدًا. وقد حققت نجاحًا كبيرًا في المقابل حين استخدمت الدقيق المطحون منزليًا مع قالب الخبز أو كخبز قالب. وإذا كنت قد وجدت طرقًا رائعة للتعامل مع الدقيق المطحون منزليًا فأرجو أن تشاركنا بها. فإمكانات استخدام الدقيق المطحون منزليًا هائلة؛ فهي ستمكّن حتى المجتمعات النائية من زراعة قمحها الخاص وصنع خبز طازج مذهل.
